\documentclass[11pt]{article}
\usepackage[T1]{fontenc}
\usepackage{lmodern}
\usepackage[margin=1in]{geometry}
\usepackage{amsmath,amssymb,amsthm}
\usepackage[hidelinks]{hyperref}
\hypersetup{pdftitle={A certified stopping point from two gradient samples},pdfauthor={},pdfsubject={AI-assisted research working paper, pending author review}}
\newtheorem{theorem}{Theorem}
\newtheorem{lemma}[theorem]{Lemma}
\newtheorem{proposition}[theorem]{Proposition}
\newcommand{\norm}[1]{\left\lVert#1\right\rVert}
\newcommand{\R}{\mathbb R}
\title{A certified stopping point from two gradient samples}
\author{}
\date{November 2026}
\begin{document}
\maketitle
\begin{abstract}
A line search can accept the value at a point while its gradient
remains above a prescribed stopping tolerance. We use the gradients
at the current and accepted trial points to certify an interior
point on the same search line. The modification makes an extra
gradient query only when that query is guaranteed to end the run.
In an exact, separated value-and-gradient model, the modification
uses at most as many requests of either type as its reference run.
An exact one-dimensional example ends in two committed steps instead
of three and uses four requests of each type in both runs. It shows
why an earlier stopping step need not save oracle requests.
\end{abstract}

\section{Reference line search and the decision to stop}

Let \(f\in C^2(\R^n)\) satisfy
\[
 \tfrac12 I\preceq\nabla^2 f(x)\preceq 4I,\qquad
 \norm{\nabla^2 f(x)-\nabla^2 f(y)}_2\le\norm{x-y}.
\]
Put \(g(x)=\nabla f(x)\). At the initial point \(x_0\) request \(f(x_0)\)
and \(g(x_0)\) and set the tolerance \(\theta=2^{-20}\norm{g(x_0)}\).
A run stops at the first \emph{committed} point where
\(\norm{g(x)}\le\theta\). The oracle replies are exact and deterministic.
A value request and a gradient request at the same point are recorded
separately, and an uncommitted trial cannot serve as a previous iterate
in the step-selection rule.

The reference \(R\) first proposes the cold step \(1/4\). At subsequent
nonterminal points it uses the most recent committed difference
\(s=x_k-x_{k-1}\) and \(y=g(x_k)-g(x_{k-1})\) to compute
\[
 b_1=\frac{s^\top s}{s^\top y},\qquad
 b_2=\frac{s^\top y}{y^\top y}.
\]
The proposal is \(b_1\) if \(b_2/b_1\ge 4/5\); otherwise it is the
minimum of the last ten computed \(b_2\) values, including the
current one. Clip the proposal to \([1/4,2]\). This fixes an
ABBmin1-type reference \cite{abb} for the comparison, including its memory.
For the Grippo--Lampariello--Lucidi test \cite{gll}, let
\(F_k=\max_{\max(0,k-9)\le j\le k}f(x_j)\).
At its current trial point \(p=x-a g(x)\), take the first step satisfying
\begin{equation}
 f(p)\le F_k-\tfrac14 a\norm{g(x)}^2.
 \label{eq:accept}
\end{equation}
It halves \(a\) after every failed value test. Rejected points
have no gradient request. At the first passing \(p\), the reference
requests \(g(p)\) and commits \(p\).

The modification \(C\) runs the same reference proposal and value
tests. It also requests \(g(p)\) at the first passing point. At the
cold step, or when \(p\) already meets the tolerance, it commits
\(p\). At other steps it may use those two gradients to certify a
point between \(x\) and \(p\). Only a certified terminal point may
replace \(p\).

\section{A certificate for a point inside the accepted segment}

\begin{lemma}\label{lem:interpolation}
Suppose the Hessian of \(f\) is \(M\)-Lipschitz in operator norm.
For \(0\le t\le 1\) and \(d=p-x\),
\begin{equation}
 \norm{g(x+td)-(1-t)g(x)-tg(p)}
 \le \tfrac12 M t(1-t)\norm{d}^2.
 \label{eq:interpolation}
\end{equation}
\end{lemma}
\begin{proof}
Let \(h(s)=g(x+s d)\), so
\(h'(s)=\nabla^2f(x+sd)d\) is Lipschitz with constant
\(M\norm{d}^2\). The interpolation error equals
\[
 (1-t)\int_0^t h'(s)\,ds
      -t\int_t^1 h'(s)\,ds .
\]
Subtract \(h'(t)\) inside both integrals. The constant terms cancel,
and bounding the two integrals by the Lipschitz estimate gives
\[
 M\norm{d}^2\left((1-t)\frac{t^2}{2}
                  +t\frac{(1-t)^2}{2}\right).
\]
This is the right side of \eqref{eq:interpolation}.
\end{proof}

Write \(g=g(x)\), \(g_p=g(p)\), and \(v=g_p-g\).
If \(v\ne0\), let \(t\) be \(-g^\top v/(v^\top v)\) clipped to
\([0,1]\). This choice minimizes \(\norm{g+t v}\) over the segment
parameter. At an endpoint value of \(t\), or when \(v=0\), keep \(p\).
For \(0<t<1\), define
\begin{equation}
 q=x-ta g,\qquad
 U=\norm{g+tv}+\tfrac12 t(1-t)a^2\norm{g}^2 .
 \label{eq:certificate}
\end{equation}
If \(U>\theta\), the modification keeps \(p\) without a request at
\(q\). If \(U\le\theta\), it constructs \(q\), requests \(g(q)\)
and checks the stopping tolerance, then requests \(f(q)\) and commits
\(q\). The supplied Hessian Lipschitz constant is one, including when
the actual constant is smaller. The bound in Lemma~\ref{lem:interpolation}
holds for every \(t\), so it applies also to this choice from measured
endpoint gradients. Consequently the extra gradient check always
succeeds in the stated model.

The new step still satisfies the value acceptance condition with step
size \(ta\). Convexity and \(f(x)\le F_k\) give
\[
 f(q)\le (1-t)f(x)+t f(p)
       \le F_k-\tfrac14 ta\norm{g}^2.
\]
This guarantee does not change the protocol's request for the actual
terminal function value. The clipping rule minimizes the norm of the
linear interpolant. The certificate adds a separate error term, and
the rule need not minimize that sum over \(t\).

\section{Stopping and request counts}

\begin{theorem}\label{thm:counts}
The reference \(R\) terminates in finitely many committed steps.
The modification \(C\) also terminates. If \(N_R,N_C\) count
committed steps and \(J_R,J_C\) count rejected value trials, let
\(T\) be one if a certified interior point is committed and zero
otherwise. Then
\[
 N_R-N_C\ge T,\qquad J_R-J_C\ge0 .
\]
Counting requests including the initial point, the differences
(reference minus modification) are
\begin{align}
 N_f^R-N_f^C&=(N_R-N_C)-T+(J_R-J_C), \label{eq:countf}\\
 N_g^R-N_g^C&=(N_R-N_C)-T. \label{eq:countg}
\end{align}
Both differences are nonnegative.
\end{theorem}

\begin{proof}
If the initial point is already terminal both runs stop with one
request of each type. Otherwise \(\theta>0\). Along a nonterminal
run, a committed step has \(s\ne0\). Strong convexity gives
\(s^\top y\ge \tfrac12\norm{s}^2>0\), so both reference proposals
are defined. For every trial
\(a\le 2(1-1/4)/4=3/8\), the smoothness estimate at \(x-a g(x)\)
implies \eqref{eq:accept}. With the clipped proposals and
halving, each line search therefore accepts a step at least
\(3/16\). Indeed a proposal passing immediately is at least \(1/4\),
whereas a final halving from a rejected trial comes from a step
larger than \(3/8\).

If the reference never stopped, every post-search value would be at
most \(F_k-c\), where \(c=(3/64)\theta^2>0\), and the window maxima
\(F_k\) could only decrease. After ten new committed values, each
value in the window is bounded by its old maximum minus \(c\).
It follows that \(F_{k+10}\le F_k-c\). Iterating this inequality
contradicts the lower bound of a strongly convex function.

For the coupled runs, suppose the committed prefixes are equal.
The proposal, the window maximum, the sequence of rejected value
tests and the first passing point \(p\) then agree. Without an
interior replacement the runs commit the same point. With one,
\(C\) stops at \(q\) while \(R\) commits the same nonterminal \(p\)
and needs at least one further step. Lemma~\ref{lem:interpolation}
rules out a failed extra gradient check. Thus there is at most one
interior replacement, and until it occurs the two committed
histories are identical. The inequalities for \(N\) and \(J\)
follow, as all rejected trials up to the replacing transaction
are included in both runs.

Each transaction requests one value and one gradient at its first
passing point, and each rejected trial costs one value request.
An interior replacement also costs one value and one gradient at
\(q\). Thus \(N_f^C=1+N_C+J_C+T\) and
\(N_g^C=1+N_C+T\), with \(T=0\) in the corresponding reference
expressions. Subtraction gives \eqref{eq:countf} and \eqref{eq:countg}.
\end{proof}

\clearpage
\section{A run with earlier stopping and equal requests}

Use the following fixed one-dimensional example. The piecewise
function and the initial point come from an earlier fixed construction.
The earlier derivation did not complete the reference comparison
beyond its second committed step. The current
calculation checks the third step separately, with the same fixed
function, initial point and tolerance. The input was selected for the
construction.
\[
 f(x)=
 \begin{cases}
 x^2/2, & x\le0,\\
 x^2/2-x^3/6, & 0\le x\le 1/2,\\
 5/48+\tfrac38(x-1/2)+\tfrac14(x-1/2)^2, & x\ge1/2 ,
 \end{cases}
 \qquad x_0=2^{-12}.
\]
The pieces agree in value and the first two derivatives at the
junctions. The second derivative is \(1\), \(1-x\) and \(1/2\)
on the three pieces. It is globally Lipschitz with constant one,
and the minimizer is \(0\).

Both runs take the cold step \(a_0=1/4\) to
\(x_1=24577/134217728\). The initial tolerance is
\(\theta=8191/35184372088832\). In one dimension the two BB
step formulas coincide whenever the previous secant is nonzero.
Between \(x_0\) and \(x_1\) the secant slope is
\(1-(x_0+x_1)/2\). Therefore the second proposal and its trial
point are
\[
 a=\frac{268435456}{268378111},\qquad
 p=-\frac{24577}{1099276742656}.
\]
The proposal lies inside the clipping interval, the first value
trial passes, and \(g(p)=p\). In particular \(p\) has not met
the stop condition. The affine interpolant of the endpoint gradients changes
sign, so its minimizing parameter is
\(t=g(x_1)/(g(x_1)-p)\in(0,1)\) and \(g(x_1)+t(g(p)-g(x_1))=0\).
The point and certificate simplify to
\begin{equation}
 q=\frac{p(g(x_1)-x_1)}{g(x_1)-p}
   =\frac{604028929}{295093865557572718592},
 \qquad U=\frac{t(1-t)a^2g(x_1)^2}{2}.
 \label{eq:explicitq}
\end{equation}
Exact rational comparison gives
\[
 \frac{g(x_1)}{\theta}>10^5,\qquad
 96<\frac{-p}{\theta}<97,\qquad
 0<\frac{U}{\theta}<\frac9{1000},\qquad
 0<\frac{g(q)}{\theta}<\frac9{1000}.
\]
Thus \(C\) commits \(q\) and stops on its second step. For the
reference, the secant between \(x_1\) and \(p\) produces
\[
 b=\frac{x_1-p}{g(x_1)-p}
   =\frac{72050996699725824}{72044400770891777}.
\]
It is also inside the clipping interval. Its first trial is
\(p-bg(p)=p(1-b)=q\), which is terminal. The remaining acceptance
checks can be written using one common value, \(f(x_0)\), because
each passing value is below it and the window still contains
\(x_0\). Direct substitution gives
\begin{align*}
 f(x_0)-\tfrac1{16}g(x_0)^2-f(x_1)&>9\cdot10^{-9},\\
 f(x_0)-\tfrac14 a g(x_1)^2-f(p)&>2\cdot10^{-8},\\
 f(x_0)-\tfrac14 b p^2-f(q)&>2\cdot10^{-8}.
\end{align*}
These comparisons also show that none of the three reference
transactions rejects its first value trial. The step table is
\[
\begin{array}{c|c|c|c|c}
  &\text{committed points after }x_0 &\text{rejected trials}
           &\text{value requests}&\text{gradient requests}\\ \hline
 R &x_1,\ p,\ q &0&4&4\\
 C &x_1,\ q     &0&4&4
\end{array}
\]
In the modified run, the requests at the uncommitted point \(p\)
account for the equality. This identity of the two runs' final
points is a property of the one-dimensional exact secants and
the verified first-trial acceptance in this example.

\section{What is being counted}

The theorem compares exact requests of each type. The work needed
to form the interpolation certificate is incurred every time it is
tested, even when no interior gradient is requested. If a value
request costs the same fixed \(\phi\) and a gradient request the
same fixed \(\gamma\) in both runs, the difference in oracle
charges is
\[
 (\phi+\gamma)(N_R-N_C-T)+\phi(J_R-J_C).
\]
Timing, memory, arithmetic, evaluation fusion and precision have
additional costs. The example has no oracle saving to offset that
extra arithmetic. The earlier trial-value rule in \cite{trial} can change the
committed history before stopping. The modification here preserves
the reference's preceding commits. Its guarantee concerns exact
separated requests. A finite-precision implementation would have to
bound the error in \(U\), the constructed point \(q\), and the
comparison with \(\theta\).

\paragraph{Status and AI disclosure.}
This AI-assisted research working paper consolidates the earlier
certified-chord derivation and completes an exact comparison for its
stored, deliberately selected test. It is pending human author review.
The companion Lean source checks the fixed rational identities,
stopping and acceptance inequalities, and arithmetic equalities for
the request counts. The execution ledger and the general coupling
are justified in the written proofs.
\begin{thebibliography}{9}
\bibitem{abb}
G. Frassoldati, L. Zanni, and G. Zanghirati,
New adaptive stepsize selections in gradient methods,
\emph{Journal of Industrial and Management Optimization},
4 (2008), pp.\ 299--312.
\href{https://doi.org/10.3934/jimo.2008.4.299}{doi:10.3934/jimo.2008.4.299}.
\bibitem{gll}
L. Grippo, F. Lampariello, and S. Lucidi,
A nonmonotone line search technique for Newton's method,
\emph{SIAM Journal on Numerical Analysis}, 23 (1986), pp.\ 707--716.
\href{https://doi.org/10.1137/0723046}{doi:10.1137/0723046}.
\bibitem{trial}
Choosing a shorter step from a line-search trial value,
anonymized working paper, 2026.
\href{https://taitai2025.github.io/assets/documents/research/trial-firsthit.pdf}
{Archived version}.
\end{thebibliography}
\end{document}
